% Section 7.2 — From Population Statistics to Spatial Information
% Combines 7.2_introduction, 7.2.1_simulated_sky_comparison, 7.2.2_frequentist_framework

\section{From Population Statistics to Spatial Information}
\label{sec:pop_to_spatial}

The previous section established that a fixed TS threshold discards a substantial population of sub-threshold sources whose aggregate properties --- including the $dN/dS$ --- are nonetheless measurable via photon-count statistics~\cite{Malyshev:2011zi, Zechlin:2015wdz}.
The $dN/dS$ is, however, a purely statistical quantity: it describes how many sources exist at each flux level but provides no information about their locations on the sky.
This section presents the conceptual framework that complements the $dN/dS$ with spatial information, transforming a one-dimensional flux distribution into a map of candidate source directions.

The idea, developed in detail in the paper body that forms the core of this chapter (Section~\ref{sec:transition_catalog}), is to use the recovered $dN/dS$ as input for generating ensembles of synthetic skies, and then to compare their pixel-wise TS distributions against the real Fermi-LAT data using a frequentist statistical test. We first describe the simulated-sky comparison procedure at a qualitative level (\S\ref{sec:simulated_sky}), and then introduce the statistical ingredients --- a per-pixel TS, a KS two-sample test, and a Quality Factor --- that make the comparison quantitative (\S\ref{sec:frequentist_framework}).
The full mathematical and computational details are presented in the paper body that follows.

\subsection{The Core Idea: Simulated-Sky Comparison}
\label{sec:simulated_sky}

As established in Section~\ref{sec:threshold_limits}, the $dN/dS$ recovered in Chapter~\ref{ch:6} extends to fluxes roughly 20 times below the Fermi-LAT detection threshold (cf.\ Section~\ref{sec:source_count}), yet it carries no spatial content whatsoever.
The question we now address is whether this population-level flux distribution can be leveraged to identify specific sky directions that are likely to host sub-threshold point sources.

The underlying idea is conceptually simple.
If the $dN/dS$ is known, one can use it as a generative model to produce realistic synthetic skies.
For each realization, source fluxes are drawn from the $dN/dS$, with the expected number of sources per flux bin given by the integral of $dN/dS$ over that bin and the actual number drawn from a Poisson distribution centered on this expectation.
Since the $dN/dS$ carries no positional information, these synthetic sources are placed at uniformly random positions on the sphere.
Each synthetic sky is then processed through the same pipeline applied to the real data: the point sources are convolved with the Fermi-LAT PSF, multiplied by the exposure map, and converted into a pixelized photon-count map on top of a modeled diffuse background (see Section~\ref{sec:transition_catalog} for full details).
By repeating this procedure 5000 times --- each time sampling a slightly different $dN/dS$ realization within the uncertainties estimated via Gaussian Process interpolation --- one obtains an ensemble of synthetic skies that collectively encode what the gamma-ray sky should look like if the recovered $dN/dS$ is correct.

The comparison between synthetic and real skies cannot be performed on a pixel-by-pixel basis, because the synthetic sources are randomly positioned and their spatial pattern never matches the real sky.
Instead, the relevant comparison is statistical: one computes a pixel-wise test statistic (TS) for every pixel in both the real and synthetic maps and then compares the resulting TS distributions.
At high TS, where bright sources dominate, the synthetic and real distributions agree well.
As one moves to lower TS values, the simulated distributions eventually diverge from the real sky, because background modeling systematics become increasingly important at faint fluxes.
The threshold at which this divergence becomes statistically significant determines how deep below the nominal catalog limit the method can reliably probe.
Pixels in the real sky whose TS exceeds this threshold are labeled ``firing pixels'' --- candidate directions that are likely to host a genuine point source.

This approach differs from earlier efforts at probabilistic gamma-ray cataloging.
The Bayesian framework PCAT~\cite{Daylan:2016tia} samples posterior distributions of full catalogs --- including source positions, fluxes, and spectra --- using transdimensional Markov Chain Monte Carlo.
While PCAT provides a complete probabilistic description of the source population in a given region of interest, its computational cost scales roughly as the square of the number of modeled sources, which limits its applicability to restricted sky areas.
In its original implementation, PCAT was applied to a $40^\circ \times 40^\circ$ patch around the North Galactic Pole, requiring approximately 250~CPU hours for a single inference run~\cite{Daylan:2016tia}. Machine learning approaches to probabilistic cataloging have also been explored~\cite{Bhat:2021wtb}, though these too have been applied to localized regions rather than to the full high-latitude sky.

The simulated-sky comparison framework adopted here takes a different route.
Rather than sampling the full catalog posterior, it reduces the problem to a comparison of one-dimensional TS distributions, which is computationally inexpensive and naturally scalable to the entire high-latitude sky.
The price for this simplification is that the method does not return individual source parameters (flux, spectrum, position uncertainty); its output is a map of candidate firing pixels at a chosen pixel resolution.
This is sufficient for the downstream application we have in mind --- cross-correlation studies (Chapter~\ref{ch:8}) --- where the relevant quantity is the number and angular distribution of candidate source directions, not the properties of individual sources.
We now turn to the statistical framework that makes this comparison quantitative.

\subsection{A Frequentist Framework: TS, KS Test, and Quality Factor}
\label{sec:frequentist_framework}

The framework outlined in the previous section is entirely frequentist.
There is no Bayesian prior on source properties in the usual sense; the $dN/dS$ acts as a generative model for the synthetic sky simulations, not as a prior to be updated.
The statistical machinery involves three ingredients: a per-pixel test statistic, a two-sample comparison test, and a meta-parameter that controls the stringency of the comparison.

The first ingredient is a per-pixel TS quantifying how much the observed photon count deviates from the background-only expectation.
For each pixel $i$, we define
\begin{equation}
	\mathrm{TS}_i = \frac{(x_i - \lambda_i)^2}{\lambda_i} \,,
	\label{eq:ts_pearson}
\end{equation}
where $x_i$ is the photon count in pixel $i$ (from either the real sky or a synthetic realization) and $\lambda_i$ is the expected count under the best-fit background-only model, consisting of a Galactic diffuse template and an isotropic component.
This definition is inspired by Pearson's $\chi^2$ statistic for Poisson-distributed counts, where the variance equals the mean, $\sigma_i^2 = \lambda_i$.
While it loosely resembles the TS used by the Fermi-LAT collaboration, there is no exact correspondence: the Fermi-LAT pipeline applies energy-dependent background renormalizations independently in each region of interest, producing a TS whose absolute value varies across the sky.
We instead use a single, globally coherent definition valid for all pixels, and we do not assume any particular probability distribution for the TS values.
The TS is treated purely as a ``signal interest label'' --- a heuristic score that flags pixels deviating from the background expectation.
Its probabilistic meaning is derived entirely from the comparison with simulated skies, not from an assumed analytic distribution.

The second ingredient is the comparison procedure itself.
As discussed in the previous section, the relevant comparison is between the overall TS distributions, not individual pixel values.
For a given minimum threshold $\mathrm{TS}^\star$, the normalized descending cumulative distributions of TS values above $\mathrm{TS}^\star$ are compared between the real sky and each synthetic realization using a two-sample Kolmogorov--Smirnov test.
The KS test determines whether the two distributions are statistically compatible at a chosen significance level $\alpha$.
At large $\mathrm{TS}^\star$, where both distributions are dominated by bright sources, they agree well.
As $\mathrm{TS}^\star$ is lowered, background systematics increasingly affect the comparison, and at some point the distributions become incompatible at the chosen confidence level.
This critical $\mathrm{TS}^\star$ defines how deep the method can probe.

The third ingredient is the Quality Factor (QF).
Because the analysis relies on an ensemble of 5000 synthetic sky realizations, not all of them will pass the KS test at a given $\mathrm{TS}^\star$ and $\alpha$.
The QF is the fraction of simulations that do pass.
A QF near unity demands that nearly all simulations agree with the data, yielding a high $\mathrm{TS}^\star$ threshold and a conservative, high-purity selection of firing pixels.
A lower QF relaxes this requirement: more simulations are allowed to fail the test, the $\mathrm{TS}^\star$ threshold drops, and the catalog extends deeper into the sub-threshold regime at the cost of including a larger fraction of spurious directions.

The pair $(\alpha, \mathrm{QF})$ thus controls the operating point of the probabilistic catalog.
Neither parameter has a uniquely correct value; both are meta-parameters that the user can tune depending on the downstream application.
Stricter choices (larger $\alpha$, higher QF) produce fewer but more reliable firing pixels.
Looser choices yield a larger statistical sample at the expense of purity.
The paper body explores the full $(\alpha, \mathrm{QF})$ plane and provides the results in digital form, so that users can select the operating point appropriate to their science case (see Section~\ref{sec:transition_catalog}).

Two further aspects of this framework deserve emphasis.
First, the TS defined above is a globally homogeneous quantity: every pixel is treated identically under the same background model, producing a TS scale whose meaning is uniform across the sky.
This contrasts with the Fermi-LAT catalog's locally computed TS, which effectively has a spatially varying significance scale.
The globally homogeneous TS is particularly advantageous for statistical analyses that combine information across large sky areas, such as the cross-correlation studies discussed in Chapter~\ref{ch:8}.
Second, at the chosen pixel resolution ($N_\mathrm{side} = 512$, corresponding to $\sim 0.12^\circ$), a single astrophysical source may produce multiple adjacent firing pixels due to PSF spreading, and conversely, multiple faint sources may contribute to a single pixel.
The method's output is therefore a map of firing pixels, not a traditional source catalog, and the relationship between firing pixels and individual sources depends on the pixel resolution and the instrument's angular resolution.

The remainder of this chapter presents the full analysis: the data selection, the detailed statistical procedure, the results for different $(\alpha, \mathrm{QF})$ operating points, and the validation against the 4FGL-DR3 catalog.
